\section{Numerical verification of the constants}
\label{sec:numerics}% Generated by verify_publication_tables.py; do not hand-edit.
\newcommand{\numBoundaryDirection}{1.224744871327}
\newcommand{\numBoundaryDirectionLimit}{1.224744871392}
\newcommand{\numBoundaryLimit}{1.2599210499}
\newcommand{\numBoundaryRatio}{1.2599210420}
\newcommand{\numCenterDistance}{0.040996138377}
\newcommand{\numCompactDirection}{0.6373774341}
\newcommand{\numCompactDirectionLimit}{0.6373774392}
\newcommand{\numCompactMultiplier}{-0.3749999965}
\newcommand{\numCorankEta}{0.715080068980}
\newcommand{\numCorankKrOne}{0.581665969555}
\newcommand{\numCorankKrTwo}{0.336668060890}
\newcommand{\numCorankSlopeFour}{-0.7150801}
\newcommand{\numHuberBoundaryRatio}{1.9999999500}
\newcommand{\numNonDiagonalCosineFive}{0.999999999713}
\newcommand{\numNonDiagonalCosineSix}{0.999999999997}
\newcommand{\numNonDiagonalDirectionFive}{0.671768177074}
\newcommand{\numNonDiagonalDirectionSix}{0.671774660292}
\newcommand{\numNonDiagonalEta}{0.644812526657}
\newcommand{\numNonDiagonalFiniteTOneFive}{0.701047395444}
\newcommand{\numNonDiagonalFiniteTOneSix}{0.701052493549}
\newcommand{\numNonDiagonalMR}{0.6717753807}
\newcommand{\numNonDiagonalSlopeFive}{-0.644805516183}
\newcommand{\numNonDiagonalSlopeSix}{-0.644811825605}
\newcommand{\numNonDiagonalTOne}{0.7010530600}
\newcommand{\numPowerThreeLimit}{1.2778862085}
\newcommand{\numPowerThreeRatio}{1.2778860451}
\newcommand{\numQuarticFour}{-0.124998347562}
\newcommand{\numQuarticOne}{0.004629629625}
\newcommand{\numQuarticThree}{0.374999996875}
\newcommand{\numQuarticTwo}{0.023437499951}
\newcommand{\numRegularLimit}{-1.7777777778}
\newcommand{\numRegularRatio}{-1.7777777766}
\newcommand{\numStrictLimit}{-0.3227486122}
\newcommand{\numStrictRatio}{-0.3227486102}


The local laws predict constants, not only exponents.  Deterministic scripts
therefore test closed-form instances at high precision; one generator
produces the displayed decimal values and machine-readable snapshots.
The full tables and execution details are in the technical supplement.  For the fixed
\(2\times2\) family
\(G=\bigl(\begin{smallmatrix}0&b_1\\b_2&c\end{smallmatrix}\bigr)\),
\(\Theta=e_1e_1^\top\), the three basic normalized limits are:
\begin{center}
\begin{tabular}{@{}lll@{}}
\toprule
Regime & predicted & computed at the stated \(\delta\)\\
\midrule
regular & \(-16/9\) & \(\numRegularRatio\), \(\delta=10^{-5}\)\\
strict corank one & \(-5/(4\sqrt{15})\) & \(\numStrictRatio\), \(\delta=10^{-8}\)\\
boundary & \(2^{1/3}\) & \(\numBoundaryRatio\), \(\delta=10^{-12}\)\\
\bottomrule
\end{tabular}
\end{center}
On the same boundary instance, standard Huber instead gives
\((1-\lambda_\delta^{\rm hub})/\delta=\numHuberBoundaryRatio\) at
\(\delta=10^{-8}\), versus the prediction \(1/\beta=2\) in
\eqref{eq:huber-boundary-rate-main}.
For the admissible profile
\(\psi'(x)=x/(1+x^3)^{1/3}\) on \(x\ge0\), the predicted \(p=3\)
constant is \(\numPowerThreeLimit\); the computed normalized displacement is
\(\numPowerThreeRatio\) at \(\delta=10^{-9}\).  This checks that the exponent
changes with the tail, not merely between the two named profiles.

For a corank-two check, take
\(A^\star=\diag(1,1,0,0)\),
\(\Theta=\diag(3/8,3/8,1,1/2)\), and \(\lambda^\star=1\).
Here \(\eta_R=\numCorankEta\) and
\(\|K_F-K_R\|_F=\numCenterDistance\).  At \(\delta=10^{-4}\), the
pseudo-Huber slope is \(\numCorankSlopeFour\), while Huber gives the water-filling
slope \(-3/5\) and recovers \(\Phi_F^\star\) to \(10^{-11}\).
Because \(P_{\rm big}(s)\) and \(\widetilde B(s)\) are constant on this
diagonal line, \(\partial_s\mathcal W=0\) proves
\(t_1=\mathcal M_R=0\).  A non-diagonal rational instance instead predicts
\(\|\mathcal M_R\|_F=\numNonDiagonalMR\), versus the scaled value
\(\numNonDiagonalDirectionSix\) at \(\delta=10^{-6}\).

The compact-root instance
\(A^\star=\diag(2,1,0)\), \(\Theta=\diag(1,-1,1)\) verifies the predicted
cubic coefficient \(-3/8\) and quadratic direction norm \(\sqrt{26}/8\).
After the second cancellation, the rational \(4\times4\) instance in the
supplement has \(\chi_2=1/3\) and converges monotonically to the quintic
coefficient \(-1/8\); its normalized quartic direction converges to the full
matrix in \eqref{eq:compact-root-quartic-main}.

The moving-family assertions cover boundary, critical, strict-side, and
oscillatory paths and verify both the cubic root and direction constants.
Complete data, matrices, and finite-scale sequences are in the supplement;
the checks test constants but imply no occurrence-frequency law.

These are CPU numerical checks, not GPU training measurements or validated
interval bounds.  The supplied verification entry point rejects optimized
Python modes that would remove assertions.  For approximate inner directions,
the relevant diagnostic is the residual of the direction actually returned,
together with the repaired primal--dual gap in the supplement; substituting a
fresh compact polar factor measures a different matrix.  Neither these local
checks nor a small tangency residual establishes an outer-training benefit.
